\section{Implementation-Ready Experimental Protocol}
\label{app:implementation-ready-protocol}

This appendix is the execution contract for E1--E9.  Values declared here are
fixed before numerical results are inspected.  A later implementation may
change a value only through a versioned protocol revision that is recorded in
the artifact manifest and applied to all affected methods.

\subsection{Identifiers, roots, and artifact layout}
\label{app:protocol-identifiers}

The protocol version is
\texttt{quark-empirical-v1}.  Its structural replicate roots are
\[
 \mathcal S_{\mathrm{struct}}
 =\{1101,1102,1103,1104,1105\}.
\]
Every root deterministically spawns the named streams
\texttt{dataset\_generation}, \texttt{dataset\_split},
\texttt{jl\_projection}, \texttt{reservoir\_parameters},
\texttt{observable\_selection}, \texttt{reset\_trajectories},
\texttt{shadow\_bases}, \texttt{measurement\_outcomes},
\texttt{classical\_baseline}, \texttt{model\_selection}, and
\texttt{replicate}.  Measurement repetitions are child seed bundles of the
structural root rather than integer offsets.

The canonical run identifier is
\[
 \texttt{<protocol>/<experiment>/<scenario>/<method>/<root>/<configuration-hash>}.
\]
The corresponding directory is
\begin{verbatim}
storage/artifacts/experiments/quark-empirical-v1/
  <E1--E9>/<scenario>/<method>/root=<root>/<configuration-hash>/
\end{verbatim}
A complete run contains the versioned feature artifact, data and split
manifests, selected readout candidate, predictions, long-form metrics, timings,
resource counters, environment, and status.  Hardware runs additionally retain
all Runtime job handles and raw shadow provenance.  Aggregate files are written
under \texttt{storage/artifacts/experiments/quark-empirical-v1/aggregate/}; they never
replace the immutable run directories.

\subsection{Common data, model, and readout objects}
\label{app:protocol-common-objects}

\paragraph{Synthetic outer split.}
For E1--E6, each process produces one ordered sequence of 6000 supervised
windows after initialization or burn-in.  Indices \(0{:}3999\),
\(4000{:}4999\), and \(5000{:}5999\) are inner train, validation, and test,
respectively.  E1 learning curves use nested prefixes of the first 5000
windows and reserve the final 20\% of each prefix for validation.  No random
shuffling is performed.

\paragraph{E1 Gaussian VARMA realization addendum.}
For the E1 reference process, the innovation covariance is \(I_3\) and the
bounded observation scale in \cref{eq:appendix-varma-observation} is
\(c_x=1\). Three Gaussian AR matrix directions are normalized to unit
spectral norm, assigned positive normalized random weights, and multiplied by
a common scalar found by a deterministic bracketed root solve so that the
companion spectral radius is \(0.7\) to numerical tolerance \(10^{-12}\).
The MA coefficients are \(B_j=0.5(0.5)^{j-1}V_j\), where the \(V_j\) are
independent Gaussian directions normalized to unit spectral norm. The
augmented ARMA transition covariance is obtained from the discrete Lyapunov
equation and the initial augmented state is drawn from the resulting stationary
Gaussian law. Thus E1 uses no burn-in approximation. Here and throughout the
controlled protocol, \(s=100\) denotes the distance between consecutive window
endpoints, not a gap added to \(w\).

The named seed list is extended append-only with
\texttt{task\_functionals}; existing stream spawn keys are unchanged. E1's
base directions \(u_3,v_3\) and decay \(\gamma\) are fixed protocol constants,
so this stream is recorded as unused for E1.

\paragraph{Reference program.}
The common reference program is
\[
 (n,R,k,\lambda,w,s)=(5,3,2,0.5,25,100),
\]
with all three reservoirs assigned the same reset rate, independent reservoir
unitary draws, ring topology, Gaussian JL projection, \(\tanh\) angle map,
and every nonidentity Pauli word of weight at most two.  The observable order is
lexicographic within weight and qubit support and is stored explicitly.  The
resulting \((N,R,K)\) tensor has \(K=105\) and flattens to \(p=315\) in
reservoir-major, observable-minor order.

\paragraph{Finite readout grid.}
For standardized features, the common grid is
\begin{equation}
 \mathcal G_{\mathrm{KRR}}
 =\{1.5,2.5,5.0\}
  \times\{0.25,0.5,1,2,4,8\}
  \times\{10^{-6},10^{-4},10^{-2},1,10^2\},
 \label{eq:empirical-krr-grid}
\end{equation}
where the entries are \((\nu,\xi,\lambda_{\mathrm K})\).  The feature scaler
and validation score are fitted on inner train and validation only.  The
lexicographically first minimizer is selected and refitted on outer train.  A
method whose representation is stochastic regenerates only the representation
stream declared by the protocol; it does not receive a larger readout search.

\paragraph{Principal controls.}
The controlled comparison uses:
\begin{enumerate}
  \item \texttt{quark-exact}: the exact NVIDIA QuaRK representation;
  \item \texttt{esn-matern}: a tanh ESN whose state dimension equals the
        reported QuaRK feature dimension, with spectral radius
        \(\{0.5,0.9\}\), input scale \(\{0.2,1.0\}\), and leak rate
        \(\{0.5,1.0\}\), followed by \cref{eq:empirical-krr-grid};
  \item \texttt{random-map-matern}: a fixed Gaussian random affine map of the
        flattened standardized window to \(p\) coordinates followed by
        coordinatewise \(\tanh\) and the same readout grid;
  \item \texttt{jl-matern}: the same \(d\to n\) projection applied at every
        time point, flattened over the window, followed by the same readout;
  \item \texttt{raw-ridge}: standardized flattened windows with ridge
        \(\alpha\in\{10^{-8},10^{-6},10^{-4},10^{-2},1,10^2,10^4\}\); and
  \item \texttt{raw-matern}: standardized flattened windows and the common
        Mat\'ern grid, only under the stated resource guardrail.
\end{enumerate}
The ESN candidate count is reported separately from the readout count.  No
method is described as compute matched unless time, memory, and candidate count
are all shown.

\subsection{Metric definitions}
\label{app:protocol-metrics}

Let \(m\) be the number of examples in the evaluated split and let
\(p=R K\).  In addition to MSE and training-normalized NRMSE from
\cref{sec:empirical-protocol}, the following quantities are retained.

\paragraph{Paired error difference.}
For method \(A\), control \(B\), and structural root \(s\),
\[
 \Delta_{A-B}^{(s)}
 =\mathrm{NRMSE}^{(s)}_A-\mathrm{NRMSE}^{(s)}_B.
\]
Negative values favor \(A\).  Pairing is valid only when data, split, and task
streams match.

\paragraph{Feature errors.}
For estimated and exact flattened features
\(\widehat\phi_i,\phi_i\in\R^p\),
\[
 \mathrm{FRMSE}_{\mathrm{coord}}
 =\sqrt{\frac1{mp}\sum_{i=1}^m\norm{\widehat\phi_i-\phi_i}_2^2},
 \qquad
 e_{2,i}=\frac{\norm{\widehat\phi_i-\phi_i}_2}{\sqrt p}.
\]
The report gives the mean, median, and 95th percentile of \(e_{2,i}\).
For a frozen readout \(h\), prediction perturbation is
\[
 \Delta_h=\frac1m\sum_{i=1}^m
 \abs{h(\widehat\phi_i)-h(\phi_i)}.
\]
Relative degradation is
\[
 D_{\mathrm{rel}}
 =\frac{\mathrm{NRMSE}_{\mathrm{estimated}}
              -\mathrm{NRMSE}_{\mathrm{exact}}}
             {\max(\mathrm{NRMSE}_{\mathrm{exact}},10^{-12})}.
\]

\paragraph{Readout diagnostics.}
For dual coefficients \(\bm\alpha\), the fitted RKHS norm is
\(\sqrt{\bm\alpha^\top K\bm\alpha}\).  Condition numbers are computed in
float64 as \(\kappa_2(K)\) and
\(\kappa_2(K+\lambda_{\mathrm K}I)\); a singular value below
\(10^{-14}\sigma_{\max}\) is reported as numerical singularity rather than
silently clipped.

\paragraph{JL diagnostic.}
The pair set \(\mathcal Q\) contains 10,000 unordered nonidentical training
point pairs generated once from the \texttt{model\_selection} stream.  In
addition to \(\widehat\Delta_{0.95}\), E3 reports the fraction whose relative
squared-distance distortion exceeds \(0.5\).  Pairs with raw squared distance
below \(10^{-12}\) are excluded and replaced deterministically.

\paragraph{Compute and resource metrics.}
Wall-clock stages are projection, feature compilation, feature execution,
feature transfer, readout selection, and final fit/prediction.  GPU timings are
synchronized.  Peak device memory is obtained from the GPU memory pool before
cleanup; host peak memory uses process resident-set high-water mark.  We also
store \(n\), physical qubits when applicable, \(p\), \(NR\) reservoir outputs,
\(NRw\) exact temporal channel updates, \(NRS\) shadow snapshots/state
preparations, and \(NRSw\) sampled temporal updates.  Hardware runs additionally
store circuit groups, submitted PUBs, total shots, completed and missing
results, transpiled depths, two-qubit gate counts, layouts, queue time, and
execution timestamps.

\subsection{Experimental master matrix}
\label{app:master-matrix}

\begin{table}[H]
\centering
\scriptsize
\caption{Implementation master matrix.  Feature runs may serve several tasks
because the same structured feature tensor is reused by multiple readouts.}
\label{tab:master-matrix}
\begin{tabularx}{\linewidth}{@{}p{0.055\linewidth}p{0.22\linewidth}p{0.25\linewidth}p{0.22\linewidth}X@{}}
\toprule
ID & Data and task combinations & Varied factors & Replication and estimator & Principal outputs \\
\midrule
E1 & Medium VARMA; future, exponential, Volterra, delays \(1,5,10,20\), cross-lag \((0,15)\) & KRR path; train prefixes \(250,500,1000,2000,5000\) & 5 roots; exact NVIDIA; one cached feature tensor/root & loss path, RKHS norm, conditioning, learning and delay curves \\
E2a & Low/high VARMA; delays \(1,10,20\) & \(\lambda=0.1,0.3,0.5,0.8,0.95\) & 5 roots; 50 exact feature runs & contraction--persistence NRMSE surface \\
E2b & VARMA/HMM/GARCH at low/medium/high persistence; shared and family-specific targets & family and persistence only & 5 roots; 45 QuaRK feature runs plus ESN/random controls & cross-family absolute and paired errors \\
E3 & Latent VARMA in sparse/distributed embeddings; exponential and Volterra & \(d=3,30,100,300,500\); \(n=3,5,8\) only at stated points & 5 roots; 90 QuaRK feature runs plus controls & error, JL distortion, time, memory \\
E4 & Seven selected E1--E3 scenarios & method only & artifact aggregation; no new features & calibrated comparison table \\
E5 & Medium VARMA; Volterra and delay 20 & one-factor \(n,R,k,\lambda\) sweeps & 5 roots; 55 exact feature runs & Pareto error, time, memory, width, dimension \\
E6 & Medium VARMA; exponential, Volterra, delay 10 & snapshots and estimator protocol & 3 structural roots; independent measurement replicates & feature and prediction degradation \\
E7 & Lorenz--63 and Mackey--Glass forecasting & horizon, observation noise, MG delay & 3 trajectory roots; exact NVIDIA and two controls & NRMSE and persistence skill \\
E8 & Beijing PM2.5, live fuel moisture, hydraulic systems & method only & 5 representation roots; exact NVIDIA and controls & dataset-level transfer table \\
E9 & \(d=n=2,w=4,R=1\) VARMA; delays 1 and 3 & execution layer only & 8 windows; \(S=192,G=6\); 2 hardware periods & feature/prediction ladder and hardware provenance \\
\bottomrule
\end{tabularx}
\end{table}

\subsection{Experiment-specific execution contracts}
\label{app:experiment-contracts}

\subsubsection{E1: \texttt{quark-e1-functional-v1}}
The feature artifact is generated once per structural root with
\texttt{ExactFeatureEstimator(complex128)} and the reference NVIDIA program.
Targets are stored independently from the feature artifact.  The regularization
path uses fixed \((\nu,\xi)=(2.5,1)\) and the eleven powers
\(10^{-8},\ldots,10^2\).  Prefix results are valid only if nested-index checks
pass.  A task with zero outer-training standard deviation is marked invalid
before fitting.  The main display is a three-panel figure: regularization path,
learning curve, and delayed-recall profile.  Full path tables and conditioning
values remain in the appendix.

\subsubsection{E2: \texttt{quark-e2-memory-v1} and
\texttt{quark-e2-family-v1}}
E2a generates separate feature tensors for every
\((\text{persistence},\lambda,\text{root})\) triple and reuses each tensor for
three delays.  E2b creates one tensor for every
\((\text{family},\text{persistence},\text{root})\) combination.  HMM states and
GARCH conditional variances are written only to teacher-target artifacts; they
never enter model inputs.  Generator stationarity checks and persistence
parameters are mandatory manifest fields.  The main display combines a
contraction--delay heat map with a cross-family paired-error panel.

\subsubsection{E3: \texttt{quark-e3-dimension-v1}}
A three-dimensional latent VARMA path is generated first.  Sparse and
distributed embeddings are then deterministic functions of the latent path and
the named nuisance/embedding streams.  Thus all dimensions for one root share
the same latent realization.  Main fixed-width runs use \(n=5\) for all five
dimensions.  Width sensitivity adds \(n=3,8\) only for \(d=100,500\); the
\(n=5\) values are reused.  Raw Mat\'ern is skipped, with status
\texttt{guardrail-not-run}, for \(d>30\).  The main display has prediction and
JL-distortion panels; detailed timing and memory tables are supplementary.

\subsubsection{E4: \texttt{quark-e4-calibration-v1}}
The aggregator fails if any expected scenario lacks at least four paired
structural roots.  Its seven rows are: E1 future, E1 exponential, E1 Volterra,
medium-HMM regime, medium-GARCH volatility, \(d=500\) sparse Volterra, and
\(d=500\) distributed Volterra.  Columns are method NRMSE mean/std, paired
difference to QuaRK, feature dimension, candidate count, feature time, and
peak memory.  Missing raw-Mat\'ern guardrail runs appear as em dashes with a
status footnote, not as failures or wins.

\subsubsection{E5: \texttt{quark-e5-resources-v1}}
The unique program list is generated by changing one field of the reference
program and deduplicating by program fingerprint.  All programs use the same
data and readout grid.  Runtime is measured after one unreported warm-up and
three timed executions of the same compiled feature plan; the median is stored.
A run that exhausts device memory is retained with its requested dimensions,
last measured allocation, and \texttt{resource-exhausted} status.  Pareto
membership uses median feature time and mean test NRMSE and is recomputed from
stored rows.

\subsubsection{E6: \texttt{quark-e6-shadows-v1}}
For the frozen-readout study, 256 test windows are selected before measurement
by equally spaced indices from the full test block.  Snapshot counts are
\(100,300,1000,3000,10000\), with ten MoM blocks, three structural roots, and
five measurement replicates.  The exact scaler and readout remain frozen.
A 16-window Aer subset, using the same bases and outcome seeds wherever
semantically possible, estimates shadow-only error from an exact channel.  The
secondary noisy-training study uses 512/128/256 train/validation/test windows,
\(S=300,1000,3000\), three structural roots, and three measurement replicates;
training and evaluation use independent measurement-outcome child streams.
Raw snapshots are retained for the Aer subset and one audit replicate per
snapshot count, not for every large GPU run.

\subsubsection{E7: \texttt{quark-e7-chaos-v1}}
After burn-in, each chaotic trajectory is standardized using the inner-training
block and bounded by a fixed \(\tanh\) map before projection.  Input windows use
stride one.  A guard interval of 20 observations separates validation from
test so that no future target crosses a split boundary.  Horizon-specific
readouts reuse the same features.  Numerical integration failure, nonfinite
state, or invariant-range escape is an integrity failure and is retained in the
ledger.  The complete matrix is in \cref{app:chaotic-protocol}.

\subsubsection{E8: \texttt{quark-e8-real-v1}}
Official train/test indices are immutable.  The final 20\% of official train is
validation.  Dataset inputs use the archive's retained window representation;
training-only standardization is applied channelwise before JL projection.
Raw Mat\'ern is attempted only if estimated float64 Gram storage is below
2 GiB and a single readout-grid pilot completes within 30 minutes.  The
aggregate requires all five paired representation roots for each stochastic
method, but no dataset is dropped because one method underperforms.  Dataset
specifications are in \cref{tab:real-world-sanity}.

\subsubsection{E9: \texttt{quark-e9-ibm-v1}}
The exact execution ladder, backend-selection score, fixed-layout rule,
Runtime-job chunking, and missing-result policy are given in
\cref{app:hardware-validation}.  The protocol submits no more than 360 unique
circuit groups and 1536 total shadow shots per hardware period.  Any provider
limit lower than the full group count triggers deterministic chunking by the
ordered group key; it does not change the structural plan.  Hardware results
are reported for every completed period, including an unsuccessful or partial
period, and are never pooled without the calibration timestamps.

\subsection{Main-paper and appendix display contract}
\label{app:display-contract}

The main paper is limited to six empirical displays after results are available:
\begin{enumerate}
  \item \textbf{Figure 3: Functional and sample-complexity profile} (E1):
        regularization, nested-sample, and delay panels;
  \item \textbf{Figure 4: Memory and dependence map} (E2): contraction--delay
        heat map and cross-family paired errors;
  \item \textbf{Figure 5: Fixed-width ambient-dimension study} (E3): NRMSE and
        JL distortion against \(d\);
  \item \textbf{Table 2: Classical calibration} (E4): the seven retained
        scenarios and transparent matching fields;
  \item \textbf{Figure 6: Resource and measurement operating points} (E5--E6):
        error--runtime Pareto panel and frozen-readout shot curve; and
  \item \textbf{Figure 7: Compact IBM transfer ladder} (E9): feature error and
        frozen-readout degradation, accompanied by one small hardware-resource
        table.
\end{enumerate}
E7 and E8 are summarized in text and shown in the appendix unless their results
change a central conclusion.

The supplementary displays include all seed-level metric tables, full KRR
paths, all E2 persistence levels, all E3 width and timing rows, complete E5
sweeps including failures, E6 shadow-only and noisy-training diagnostics,
chaos horizon/noise panels, dataset-level E8 predictions, and the complete E9
job, calibration, layout, depth, count, and missing-result ledger.  Negative
and failed runs are retained rather than hidden.

\subsection{Run completion and audit rules}
\label{app:completion-rules}

A scenario is complete only when all expected manifests pass checksum and
ordering validation, all selected readouts can be reproduced from stored
features, and every displayed aggregate maps to immutable run identifiers.
Resumable hardware jobs are incomplete until decoded results or a terminal
provider status are stored.  When fewer than the required replicates finish,
the result is shown as incomplete and no mean-only replacement is reported.
All plotting and table scripts consume long-form aggregate files and are
forbidden from reading ad hoc notebook variables.
